% !TeX root = ../navier-stokes-manuscript.tex
% ============================================================
% 2.4 The Spatial Column Inside Critical Height
% ============================================================

\subsection{The Spatial Column Inside Critical Height}\label{sub:TheSpatialColumnInsideCriticalHeight}

To follow the viscous link from a temporal response into space, take the scalar
energy at an arbitrary spatial height \(s\ge0\):
\[
  E_s(t)
  :=
  \frac12\norm{\Lambda^su(t)}_2^2,
\]
In particular,
\[
  H(t)=E_{1/2}(t).
\]
Keeping \(s\) free lets the same calculation be reused after the critical row.

Apply \(\Lambda^s\) to the full momentum equation and pair with
\(\Lambda^su\).  The temporal contribution is
\[
  \pair{\Lambda^su}{\Lambda^s\partial_tu}
  =
  \frac{d}{dt}\frac12\norm{\Lambda^su}_2^2
  =
  E_s'.
\]
Viscosity commutes with \(\Lambda^s\), and spatial decay gives
\[
  \nu\pair{\Lambda^su}{\Lambda^s\Delta u}
  =
  -\nu\norm{\Lambda^{s+1}u}_2^2
  =
  -2\nu E_{s+1}.
\]
Pressure cancels in this global scalar balance only after integration by parts
uses incompressibility:
\[
  \pair{\Lambda^su}{\Lambda^s\nabla p}
  =
  -\pair{\nabla\cdot\Lambda^su}{\Lambda^sp}
  =
  0.
\]
It remains present in the local equation and in every mixed coordinate.
Accordingly, keep the signed transfer in its pressure-complete form:
\[
  \mathcal P_s(t)
  :=
  -\pair{\Lambda^su}
  {\Lambda^s\bigl((u\cdot\nabla)u+\nabla p\bigr)}.
\]
The resulting identity is
\[
  E_s'
  =
  \mathcal P_s-2\nu E_{s+1}.
\]
One derivative in time has reached one full Sobolev level higher in space, but
only together with the signed transfer at the original height.

At \(s=1/2\), this becomes
\[
  H'
  =
  \mathcal P_{1/2}-2\nu E_{3/2}.
\]
The scale-invariant height now touches \(E_{3/2}\).  Its mere appearance is not
a bound: transfer and dissipation can cancel inside \(H'\).  Differentiate the
whole row, not the isolated dissipative term.

The next energy row is
\[
  E_{3/2}'
  =
  \mathcal P_{3/2}-2\nu E_{5/2},
\]
and substitution into the derivative of the critical row yields
\begin{align*}
  H''
  &=
  \mathcal P_{1/2}'-2\nu E_{3/2}'
  \\
  &=
  \mathcal P_{1/2}'
  -2\nu\mathcal P_{3/2}
  +(2\nu)^2E_{5/2}.
\end{align*}
The response \(H''\) reaches \(E_{5/2}\), while retaining the transfer from
each preceding height.  One more derivative uses
\[
  E_{5/2}'
  =
  \mathcal P_{5/2}-2\nu E_{7/2}
\]
and gives
\[
  H'''
  =
  \mathcal P_{1/2}''
  -2\nu\mathcal P_{3/2}'
  +(2\nu)^2\mathcal P_{5/2}
  -(2\nu)^3E_{7/2}.
\]
The fourth response continues the same recurrence:
\[
  H^{(4)}
  =
  \mathcal P_{1/2}'''
  -2\nu\mathcal P_{3/2}''
  +(2\nu)^2\mathcal P_{5/2}'
  -(2\nu)^3\mathcal P_{7/2}
  +(2\nu)^4E_{9/2}.
\]
At each step, the highest spatial energy rises once and the entire lower
transfer history remains in the row.

This produces a triangular identity.
\begin{proposition}[Critical-height spatial ascent]\label[proposition]{prop:CriticalHeightSpatialAscent}
For every integer \(n\ge1\),
\[
  H^{(n)}
  =
  (-2\nu)^nE_{n+1/2}
  +
  \sum_{j=0}^{n-1}
  (-2\nu)^j
  \partial_t^{\,n-1-j}\mathcal P_{j+1/2}.
\]
\end{proposition}
\begin{proof}
The critical balance is the case \(n=1\).  If the formula holds at order
\(n\), differentiate it and use
\[
  E_{n+1/2}'
  =
  \mathcal P_{n+1/2}-2\nu E_{n+3/2}.
\]
The new transfer extends the sum by one term, while the dissipative part becomes
\((-2\nu)^{n+1}E_{n+3/2}\).  This is the formula at order \(n+1\).
\end{proof}
The proposition is an identity, not an estimate.  Knowing \(H^{(n)}\) cannot
isolate \(E_{n+1/2}\) until the complete signed transfer sum is controlled.

Critical scaling was homogeneous; continuation uses inhomogeneous Sobolev
spaces.  For \(s>0\), the two are joined by
\[
  \norm{u(t)}_{H^s}^2
  \simeq
  \norm{u(t)}_2^2+\norm{\Lambda^su(t)}_2^2
  =
  \norm{u(t)}_2^2+2E_s(t).
\]
The kinetic-energy identity pays the low-frequency term:
\[
  \frac12\norm{u(t)}_2^2
  +
  \nu\int_0^t\norm{\nabla u(\tau)}_2^2\,d\tau
  =
  \frac12\norm{u_0}_2^2.
\]
Hence a bound for \(E_s(t)\) completes an \(H^s\) bound, while time-integrability
of \(E_s\) completes the high-frequency part of \(L^2_tH^s_x\) on a finite
interval.

The ascent begins at \(1/2\) because scaling selected it, and the first response
reaches \(3/2\).  At the terminal face the logic reverses: whole-terminal
Sobolev rows first bound these two heights, from which \(H\) and its first
dissipative level follow.

Each finite number of responses reaches only a finite spatial height.  To keep
every finite coordinate of the derivative family, the same velocity must
occupy every finite temporal and spatial Sobolev row.
